% 1-2-problem.tex
%  Budget: 1pp.
%  Rules: every numeral in \lr{}; cite only from references.bib;
%  one unit at a time -- see thesis/WRITING_HANDOFF.md and APPROVALS.md.

\section{بیان مسئله}

مسئله‌ی محوری این پژوهش، ارائه‌ی پیش‌بینی دقیق و قابل‌اتکا از قیمت روز-پیش برق در
بازار آلمان است؛ آن هم بر روی دو افق هدف: برداری از \lr{24} قیمت ساعتی برای روز
پیش‌ِرو، و متوسط بار پایه‌ی روزانه. در حالت آرمانی، بازیگران بازار باید بتوانند
در همه‌ی شرایط بازار \lr{—} از دوره‌های آرام تا دوره‌های پرتنش \lr{—} به
پیش‌بینی‌هایی دست یابند که هم دقیق باشند و هم شواهدی شفاف درباره‌ی این‌که کدام
روش و چرا بهتر عمل می‌کند، در اختیار بگذارند.

اما آنچه دستیابی به این هدف را دشوار می‌سازد، ماهیت خاص سری قیمت برق است.
نوسان‌های شدید، جهش‌های ناگهانی و جابه‌جایی میان رژیم‌های آرام و پرتنش سبب
می‌شوند که عملکرد یک مدل واحد و ثابت در تمام شرایط بازار پایدار نماند
\cite{jiang_review_2018}. افزون بر این، رشد سریع روش‌های پیش‌بینی با ضعف‌هایی
روش‌شناختی در ارزیابی همراه بوده و تشخیص این‌که کدام روش واقعاً برتر است را دشوار
ساخته است \cite{lago_forecasting_2021}. مرورهای اخیر نیز کیفیت داده،
تفسیرپذیری مدل‌ها و تعمیم‌پذیری آن‌ها به شرایط تازه را از جمله چالش‌های همچنان
گشوده‌ی این حوزه برمی‌شمارند \cite{oconnor_review_2025}. پیامد این کاستی‌ها آن
است که خطای پیش‌بینی مستقیماً به تصمیم‌های نادرست و زیان مالی می‌انجامد، انتخاب
مدل بر پایه‌ای نااستوار صورت می‌گیرد، و مدل‌هایی که در یک دوره عملکرد خوبی
داشته‌اند ممکن است با تغییر شرایط بازار به‌طور ناگهانی و بی‌هشدار دچار افت شوند.

درست به همین سبب، یک مطالعه‌ی تک‌مدلی که تنها بر دقت در همان توزیع داده‌ی آموزش
تمرکز کند برای پاسخ‌گویی به این مسئله کافی نیست. نخست، از آنجا که برتریِ نسبیِ
خانواده‌های گوناگون مدل \lr{—} از روش‌های آماری کلاسیک تا یادگیری ماشین و یادگیری
عمیق \lr{—} به شرایط بازار و افق پیش‌بینی وابسته است
\cite{oconnor_review_2025, jiang_review_2018}، ترکیب مدل‌ها و تطبیق آن‌ها با
رژیم‌های متفاوت بازار می‌تواند از هر مدل منفرد راه‌گشاتر باشد. دوم، دقت بالا
به‌تنهایی بسنده نیست؛ در یک تصمیم مالیِ پرمخاطره، دانستنِ این‌که کدام عوامل
پیش‌بینی را شکل می‌دهند به‌اندازه‌ی خودِ پیش‌بینی اهمیت دارد. سوم، عملکرد خوب در
گذشته تضمینی برای پایداری در آینده نیست، و رفتار مدل هنگام تغییر ساختاری بازار
باید به‌صراحت سنجیده شود.

با توجه به این ملاحظات، مسئله‌ی اصلی این پژوهش را می‌توان چنین صورت‌بندی کرد: توسعه و
ارزیابی نظام‌مند مجموعه‌ای از مدل‌های پیش‌بینی قیمت روز-پیش برای بازار برق آلمان،
در چارچوب یک پروتکل ارزیابی استاندارد، تکرارپذیر و متکی بر بنچمارک عمومی. این
مسئله چند بُعد به‌هم‌پیوسته دارد: دقت پیش‌بینی در مقایسه با روش‌های مرجع
منتشرشده در ادبیات؛ شیوه‌ی پیش‌بینی بار پایه‌ی روزانه، به‌صورت مستقیم یا از راه
تجمیع پیش‌بینی‌های ساعتی؛ رفتار مدل‌ها در رژیم‌های متفاوت بازار؛ پایداری آن‌ها در
برابر تغییر توزیع داده در گذر زمان؛ و تفسیرپذیری عوامل مؤثر بر پیش‌بینی. این ابعاد
در قالب سؤالات پژوهش، در بخش بعد به‌صورت رسمی صورت‌بندی می‌شوند.
